% Source: Wald GR, physical PDF p. 210
%         (printed p. 197).
% Coverage: A spacetime that violates strong causality.
% Figures/tables/footnotes: Figure 8.8; no tables or footnotes.
% Status: source-reviewed and compile-clean.
% Uncertainties: none.

\begingroup
\renewcommand{\thefigure}{8.8}
\begin{figure}[H]
  \centering
  \begin{tikzpicture}[x=.82cm,y=.82cm,>=Latex,line cap=round,font=\small]
    \draw[thick] (-1.50,4.12) arc[start angle=180,end angle=360,x radius=1.50,y radius=.32];
    \draw[thick] (-1.50,4.12) arc[start angle=180,end angle=0,x radius=1.50,y radius=.32];
    \draw[thick] (-1.50,.34) arc[start angle=180,end angle=360,x radius=1.50,y radius=.32];
    \draw[dashed] (-1.50,.34) arc[start angle=180,end angle=0,x radius=1.50,y radius=.32];
    \draw[thick] (-1.50,.34) -- (-1.50,4.12);
    \draw[thick] (1.50,.34) -- (1.50,4.12);
    \draw[dashed] (-1.50,2.40) arc[start angle=180,end angle=0,x radius=1.50,y radius=.27];
    \draw[thick] (-1.50,2.40) arc[start angle=180,end angle=360,x radius=1.50,y radius=.27];
    \foreach \x/\y/\a in {-.85/3.30/0,0/3.30/0,.85/3.30/0,-.85/2.72/90,0/2.72/90,.85/2.72/90,-.85/1.70/90,0/1.70/90,.85/1.70/90,-.85/1.12/0,0/1.12/0,.85/1.12/0} {
      \begin{scope}[shift={(\x,\y)},rotate=\a,scale=.70]
        \draw (-.22,-.35) -- (0,0) -- (-.22,.35);
        \draw (.22,-.35) -- (0,0) -- (.22,.35);
        \draw (-.22,-.35) .. controls (-.04,-.39) and (.04,-.39) .. (.22,-.35);
        \draw (-.22,.35) .. controls (-.04,.39) and (.04,.39) .. (.22,.35);
      \end{scope}
    }
    \draw[thick] (-1.13,2.22) .. controls (-.54,2.23) and (.14,1.92) .. (1.08,2.18);
    \fill (-.37,2.13) circle (1.55pt) node[below] {\(p\)};
    \draw[fill=white] (.36,2.03) circle (2.0pt);
    \draw[->] (2.08,2.02) -- (.47,2.00);
    \node[anchor=west] at (2.08,1.82) {remove point};
  \end{tikzpicture}
  \caption{A spacetime which violates strong causality. The light cones \lq\lq tip
  over\rq\rq\ sufficiently that the curve drawn through \(p\) is a null geodesic
  (which, however, is not a closed curve because of the point removed from the
  manifold). There exist no closed causal curves in this spacetime, but there are
  causal curves through \(p\) which come \lq\lq arbitrarily close\rq\rq\ to intersecting
  themselves.}
  \label{fig:8.8}
\end{figure}
\endgroup
